\documentclass[10pt,a4paper]{amsart}
\usepackage[margin=19mm]{geometry}
\usepackage{amsmath,amssymb,mathtools}
\usepackage[scr=rsfs,cal=euler]{mathalfa}
\usepackage{xfrac}
\usepackage[unicode,hidelinks,bookmarks=false]{hyperref}
\newcommand{\ie}{i.e.\ }
\newcommand{\cf}{cf.\ }
\newcommand{\resp}{respectively\ }
\newcommand{\Subsection}{\subsection}
\providecommand{\nobreakdash}{\nobreak\hbox{-}}
\setlength{\emergencystretch}{2em}
\multlinegap0pt
\usepackage{amssymb}
\usepackage{mathtools}
\newtheorem{proposition}{Proposition}
\newtheorem{definition}{Definition}
\newtheorem{lemma}{Lemma}
\newtheorem{theorem}{Theorem}
\newcommand{\Hm}{\mathbb H_m}
\newcommand{\Pm}{\mathbb P_m}
\newcommand{\Tr}{\operatorname{Tr}}
\newcommand{\altPhi}{\widehat{\Phi}}
\title{Alternative-mean trace divergences: geometry, data processing, and barycenters}
\author{Trung Dung Vuong, Hiroki Shudo,, Hiroyuki Osaka}
\date{Source: arXiv:2609.00034v1 (CC BY 4.0)}
\begin{document}
\maketitle

\noindent\emph{Source and licence.} Alternative-mean trace divergences: geometry, data processing, and barycenters. Version arXiv:2609.00034v1 (CC BY 4.0), distributed by arXiv under the Creative Commons Attribution 4.0 International licence, http://creativecommons.org/licenses/by/4.0/, as recorded in the arXiv licence metadata for this exact version and shown on the abstract page https://arxiv.org/abs/2609.00034v1. Full licence notice: \texttt{licenses/arxiv-2609.00034-CC-BY-4.0.txt}.\par\smallskip

\noindent\textit{Editorial scope.} Counted unit: the article's theorem thm:certificate (source0.tex lines 2309--2319). The article's complete original proof of it is delivered with it (source0.tex lines 2321--2597, one \texttt{proof} environment of 277 lines). No mathematics is written, repaired or re-derived: the statement and the proof are the article's own lines, in the article's own notation, and no retained line is substituted or re-flowed.  Retained uncounted as the source-local context the proof needs: lines 834--838; lines 1731--1739; lines 1772--1784; lines 1946--1952; lines 1954--1971; lines 1963--1965; lines 2114--2128; lines 2133--2141; lines 2152--2169.  Nothing was omitted from any retained span, so the package carries no omission record.  No retained line contains a citation, so the wrapper carries no bibliography and no bibliography entry.  No retained line contains a figure environment, an image-inclusion command or drawing code, so no image is delivered and the package declares no asset dependency.  Editorial formatting only, in addition: an AMS \texttt{amsart} preamble that restates the article's own declarations and macros used by the retained lines, namely the environments \texttt{proposition}, \texttt{definition}, \texttt{lemma}, \texttt{theorem} on separate counters, together with the standard packages those lines need.  The retained environments print in this document as Proposition 1, Definition 1, Lemma 1, Theorem 1 in order of appearance, whereas the article prints the same items with the numbering of its own full document.  The complete native source of the article as distributed in the arXiv e-print for this exact version is bundled unchanged as \texttt{sources/208995/source0.tex}.
	\begin{equation}\label{eq:representing-measure}
		f(x)=\int_{[0,1]}h_\lambda(x)\,d\mu_f(\lambda),
		\qquad
		s=f'(1)=\int_{[0,1]}\lambda\,d\mu_f(\lambda).
	\end{equation}

	\begin{align}
		H_{f,A}(X)
		&:=\Tr\!\left(Aq(R_A(X))\right), \label{eq:HfA}\\
		H_{f,\omega}(X)
		&:=\sum_{j=1}^n w_jH_{f,A_j}(X), \label{eq:Hfomega}\\
		\mathcal E_f(X)
		&:=\sum_{j=1}^n w_j\altPhi_f(A_j,X) \notag\\
		&=C_\omega+s\Tr X-H_{f,\omega}(X), \label{eq:energy}
	\end{align}

\begin{proposition}
		\label{prop:closed-cone}
		For every nontrivial normalized operator monotone function $f$, the set
		\[
		\operatorname*{argmin}_{X\succeq0}\mathcal E_f(X)
		\]
		is nonempty and compact.  Moreover,
		\begin{equation}\label{eq:open-closed}
			\inf_{X\in\Pm}\mathcal E_f(X)
			=\min_{X\succeq0}\mathcal E_f(X),
		\end{equation}
		and $X=0$ is not a minimizer.
	\end{proposition}

	\begin{equation}\label{eq:Kg}
		\mathcal K_g(A;X)
		:=A^{1/2}\mathcal S_{C_A(X)}\!\left[
		A^{-1/2}Dg(R_A(X))[A]A^{-1/2}
		\right]A^{1/2},
		\qquad X\in\Pm.
	\end{equation}

 \begin{proposition}
		\label{prop:euler}
		Under the standing assumptions, recall that $q(x)=f(x)^2$ and
		$s=f'(1)\in(0,1)$.  For $A,X\in\Pm$ and $Y\in\Hm$,
		\begin{equation}\label{eq:first-variation}
			D_X\altPhi_f(A,X)[Y]
			=\Tr\!\left((sI-\mathcal K_q(A;X))Y\right).
		\end{equation}
		Hence every local minimizer $X_*\in\Pm$ satisfies the barycenter equation
		\begin{equation}\label{eq:euler}
			sI=\sum_{j=1}^n w_j\mathcal K_q(A_j;X_*).
		\end{equation}
		Conversely, if $\mathcal E_f$ is convex on $\Pm$, every solution of
		\eqref{eq:euler} is a global minimizer over $\Pm$.  If $\mathcal E_f$ is
		strictly convex on $\Pm$, then \eqref{eq:euler} has at most one solution;
		whenever such a solution exists, it is the unique global minimizer over
		$\Pm$.
\end{proposition}

		\begin{equation}
			sI=\sum_{j=1}^n w_j\mathcal K_q(A_j;X_*).
		\end{equation}

\begin{definition}\label{def:certificate}
		Let $m\ge1$ be an integer.  We say that $f$ satisfies the
		\emph{order-$m$ residual certificate} if there is a real constant
		$c\in\mathbb R$ such that
		\begin{equation}\label{eq:residual}
			G(x):=f(x)^2-cx^2
		\end{equation}
		is nonnegative, nonconstant, and matrix concave of order $m$ on
		$(0,\infty)$.  If, for the same constant $c$, the residual $G$ is a
		nonconstant operator monotone function on $(0,\infty)$, we say that $f$
		satisfies the \emph{dimension-free residual certificate}.  Here the word
		``certificate'' means that the stated property of the explicitly
		checkable residual $G$ will serve below as a sufficient condition for
		uniqueness and positive definiteness of the barycenter.
\end{definition}

\begin{lemma}
	\label{lem:positive-matrix-concavity}
	Let $h:(0,\infty)\to[0,\infty)$ be matrix concave of order $m$.
	Then $h$ is matrix monotone of order $m$.  If, in addition, $h$
	is real analytic and nonconstant, then
	\[
	h'(x)>0,\qquad x>0.
	\]
\end{lemma}

	\begin{lemma}
		\label{lem:certificate-constant}
	 If $f$ satisfies the order-$m$ residual certificate, then the real
		constant $c$ in \eqref{eq:residual} is uniquely determined by $f$, and
		\begin{equation}\label{eq:canonical-c}
			c=\left(\lim_{x\to\infty}\frac{f(x)}{x}\right)^2
			=\mu_f(\{1\})^2,
			\qquad 0\le c<s<1,
		\end{equation}
		where $\mu_f$ is the representing measure in
		\eqref{eq:representing-measure}.  Moreover,
		\begin{equation}\label{eq:normalized-residual}
			\widetilde G(x):=\frac{G(x)}{1-c}
		\end{equation}
		is normalized, nonnegative, nonconstant, matrix concave of order $m$, and
		matrix monotone of order $m$.  Under the dimension-free certificate,
		$\widetilde G$ is operator monotone.
	\end{lemma}

\begin{theorem}
	\label{thm:certificate}
	Let $m\ge1$, and suppose that $f$ satisfies the order-$m$ residual
	certificate of Definition~\ref{def:certificate}.  Then, for every
	integer $n\ge1$, every positive probability vector
	$\omega=(w_1,\ldots,w_n)$, and every
	$A_1,\ldots,A_n\in\Pm$, the continuous extension of
	$\mathcal E_f$ to $\overline{\mathbb{P}}_m$ has a unique global minimizer
	$X_*$.  Moreover, $X_*\in\Pm$ and is the unique solution of the
	barycenter equation \eqref{eq:euler}.
\end{theorem}

 \begin{proof}
		Fix $A\in\Pm$.  Since
		\[
		q(x)=f(x)^2=cx^2+G(x),
		\]
		functional calculus, cyclicity of the trace, and the Riccati identity
		$R_A(X)AR_A(X)=X$ give
		\begin{align}
			H_{f,A}(X)
			&=c\Tr\!\left(AR_A(X)^2\right)
			+\Tr\!\left(AG(R_A(X))\right) \notag\\
			&=c\Tr\!\left(R_A(X)AR_A(X)\right)
			+\Tr\!\left(AG(R_A(X))\right) \notag\\
			&=c\Tr X+\Tr\!\left(AG(R_A(X))\right).
			\label{eq:H-residual}
		\end{align}
		
		We first prove strict concavity.  Let $X,Y\in\Pm$ with $X\ne Y$,
		let $0<\theta<1$, and set
		\[
		X_\theta=(1-\theta)X+\theta Y,
		\qquad
		\overline R_\theta
		=(1-\theta)R_A(X)+\theta R_A(Y).
		\]
		Since $C_A$ is linear and the square-root map is operator concave,
		\begin{equation}\label{eq:strict-RA-concavity}
			R_A(X_\theta)\succeq\overline R_\theta.
		\end{equation}
		This inequality is strict in the sense that its two sides are not
		equal.  Indeed, put
		\[
		P=C_A(X)^{1/2},
		\qquad
		Q=C_A(Y)^{1/2}.
		\]
		Equality in \eqref{eq:strict-RA-concavity} would imply
		\[
		C_A(X_\theta)^{1/2}
		=(1-\theta)P+\theta Q.
		\]
		Squaring both sides and using
		\[
		C_A(X_\theta)
		=(1-\theta)P^2+\theta Q^2
		\]
		would give
		\[
		\theta(1-\theta)(P-Q)^2=0.
		\]
		Hence $P=Q$, so $C_A(X)=C_A(Y)$ and therefore $X=Y$, a
		contradiction.
		
		By Definition~\ref{def:certificate} and
		Lemma~\ref{lem:positive-matrix-concavity}, the residual $G$ is both
		matrix concave and matrix monotone of order $m$.  Therefore,
		\[
		G(R_A(X_\theta))
		\succeq G(\overline R_\theta)
		\succeq
		(1-\theta)G(R_A(X))+\theta G(R_A(Y)).
		\]
		Set
		\[
		D_1
		:=G(R_A(X_\theta))-G(\overline R_\theta)
		\]
		and
		\[
		D_2
		:=G(\overline R_\theta)
		-(1-\theta)G(R_A(X))-\theta G(R_A(Y)).
		\]
		Then $D_1,D_2\succeq0$.  Moreover, $D_1\ne0$.  Otherwise,
		\[
		G(R_A(X_\theta))=G(\overline R_\theta).
		\]
		By Lemma~\ref{lem:certificate-constant}, $G'(x)>0$ for every $x>0$,
		so $G$ is injective and has an inverse on its range.  Applying this
		inverse by functional calculus would yield
		\[
		R_A(X_\theta)=\overline R_\theta,
		\]
		contrary to the strictness of \eqref{eq:strict-RA-concavity}.
		
		It follows that $D:=D_1+D_2$ is positive semidefinite and nonzero.
		Since $A\succ0$,
		\[
		\Tr(AD)
		=\Tr\!\left(A^{1/2}DA^{1/2}\right)>0.
		\]
		Consequently,
		\[
		\Tr\!\left(AG(R_A(X_\theta))\right)
		>
		(1-\theta)\Tr\!\left(AG(R_A(X))\right)
		+\theta\Tr\!\left(AG(R_A(Y))\right).
		\]
		Thus
		\[
		X\longmapsto\Tr\!\left(AG(R_A(X))\right)
		\]
		is strictly concave.  Since $c\Tr X$ is affine,
		\eqref{eq:H-residual} shows that $H_{f,A}$ is strictly concave.
		Hence the positively weighted sum $H_{f,\omega}$ is strictly concave,
		and \eqref{eq:energy} shows that $\mathcal E_f$ is strictly convex
		on $\Pm$.
		
		We next show that no minimizer can lie on the boundary of the positive
		cone.  Let
		\[
		X_0\in\overline{\mathbb{P}}_m\setminus\Pm,
		\]
		so that $X_0\succeq0$ is noninvertible, and for $\varepsilon>0$ set
		\[
		X_\varepsilon=X_0+\varepsilon I,
		\qquad
		C_\varepsilon=C_A(X_\varepsilon),
		\qquad
		R_\varepsilon=R_A(X_\varepsilon).
		\]
		Fix $0<\varepsilon_0<1$.  The map
		\[
		\varepsilon\longmapsto R_A(X_0+\varepsilon I)
		\]
		is continuous on $[0,\varepsilon_0]$.  Its image is therefore
		bounded, so there exists $\kappa>0$ such that
		\[
		R_\varepsilon\preceq\kappa I,
		\qquad 0<\varepsilon\le\varepsilon_0.
		\]
		Put
		\[
		a:=\lambda_{\min}(A),
		\qquad
		b:=\lambda_{\max}(A),
		\qquad
		d_\kappa:=G'(\kappa)>0.
		\]
		
		Matrix monotonicity of order $m$ implies that
		$DG(R_\varepsilon)$ is a positive linear map on $\mathbb H_m$.
		Indeed, if $Z\succeq0$ and $t>0$, then
		\[
		R_\varepsilon\preceq R_\varepsilon+tZ,
		\]
		so
		\[
		\frac{G(R_\varepsilon+tZ)-G(R_\varepsilon)}{t}\succeq0.
		\]
		Letting $t\downarrow0$ and using the closedness of the positive
		semidefinite cone gives
		\[
		DG(R_\varepsilon)[Z]\succeq0.
		\]
		In particular, since $A\succeq aI$,
		\[
		DG(R_\varepsilon)[A]
		\succeq aDG(R_\varepsilon)[I]
		=aG'(R_\varepsilon).
		\]
		Scalar concavity makes $G'$ nonincreasing.  Since
		$\sigma(R_\varepsilon)\subset(0,\kappa]$, functional calculus gives
		\[
		G'(R_\varepsilon)\succeq G'(\kappa)I=d_\kappa I.
		\]
		Using also $A^{-1}\succeq b^{-1}I$, we obtain
		\begin{align}
			A^{-1/2}DG(R_\varepsilon)[A]A^{-1/2}
			&\succeq
			aA^{-1/2}G'(R_\varepsilon)A^{-1/2} \notag\\
			&\succeq
			ad_\kappa A^{-1} \notag\\
			&\succeq
			\frac{a}{b}d_\kappa I.
			\label{eq:DG-lower-bound}
		\end{align}
		
		By \eqref{eq:Kg}, positivity of
		$\mathcal S_{C_\varepsilon}$ and
		\[
		\mathcal S_{C_\varepsilon}[I]
		=\frac12C_\varepsilon^{-1/2}
		\]
		now yield
		\[
		\mathcal K_G(A;X_\varepsilon)
		\succeq
		\frac{ad_\kappa}{2b}
		A^{1/2}C_\varepsilon^{-1/2}A^{1/2}.
		\]
		Taking traces and using $A\succeq aI$, we obtain
		\begin{align}
			\Tr\mathcal K_G(A;X_\varepsilon)
			&\ge
			\frac{ad_\kappa}{2b}
			\Tr\!\left(AC_\varepsilon^{-1/2}\right) \notag\\
			&\ge
			\frac{a^2d_\kappa}{2b}
			\Tr C_\varepsilon^{-1/2}.
			\label{eq:KG-boundary-blowup}
		\end{align}
		
		Since $A\succ0$, the congruence
		\[
		C_A(X_0)=A^{1/2}X_0A^{1/2}
		\]
		is singular whenever $X_0$ is singular.  Moreover,
		\[
		C_\varepsilon=C_A(X_0)+\varepsilon A
		\longrightarrow C_A(X_0).
		\]
		Hence
		\[
		\lambda_{\min}(C_\varepsilon)\longrightarrow0,
		\]
		and therefore
		\[
		\Tr C_\varepsilon^{-1/2}
		\ge
		\lambda_{\min}(C_\varepsilon)^{-1/2}
		\longrightarrow+\infty.
		\]
		It follows from \eqref{eq:KG-boundary-blowup} that
		\[
		\Tr\mathcal K_G(A;X_\varepsilon)
		\longrightarrow+\infty.
		\]
		
		Differentiating \eqref{eq:H-residual} in the direction $I$ gives
		\[
		DH_{f,A}(X_\varepsilon)[I]
		=mc+\Tr\mathcal K_G(A;X_\varepsilon)
		\longrightarrow+\infty.
		\]
		Applying this argument to each $A_j$, summing with the positive
		weights, and using \eqref{eq:energy}, we obtain
		\[
		\frac{d}{d\varepsilon}
		\mathcal E_f(X_0+\varepsilon I)
		=
		sm-DH_{f,\omega}(X_0+\varepsilon I)[I]
		\longrightarrow-\infty.
		\]
		Thus there exists $\delta>0$ such that
		\[
		\frac{d}{d\varepsilon}
		\mathcal E_f(X_0+\varepsilon I)\le-1,
		\qquad 0<\varepsilon<\delta.
		\]
		Define
		\[
		\varphi(t):=\mathcal E_f(X_0+tI).
		\]
		For $0<\eta<\varepsilon<\delta$,
		\[
		\varphi(\varepsilon)-\varphi(\eta)
		=\int_\eta^\varepsilon\varphi'(t)\,dt
		\le-(\varepsilon-\eta).
		\]
		Letting $\eta\downarrow0$ and using continuity of the closed-cone
		extension gives
		\[
		\mathcal E_f(X_0+\varepsilon I)
		-\mathcal E_f(X_0)
		\le-\varepsilon<0.
		\]
		Hence no matrix in
		$\overline{\mathbb{P}}_m\setminus\Pm$ can minimize $\mathcal E_f$.
		
		Proposition~\ref{prop:closed-cone} supplies a global minimizer in
		$\overline{\mathbb{P}}_m$, and the preceding argument shows that it belongs to
		$\Pm$.  Strict convexity makes this minimizer unique.
		Proposition~\ref{prop:euler} then shows that it satisfies
		\eqref{eq:euler}; conversely, every solution of \eqref{eq:euler} is
		this unique minimizer.
\end{proof}

\end{document}
